\documentstyle[11pt]{article}

\renewcommand{\thefootnote}{\fnsymbol{footnote}}
\renewcommand{\baselinestretch}{1.0}
\newenvironment{comment}{\setbox0=\vbox\bgroup}{\egroup}
\newcommand{\con}{\rightarrow}
\newcommand{\cd}{\mbox{$\stackrel{d}{\rightarrow}$}}
\newcommand{\ep}{\varepsilon}
\newenvironment{singlespace}{\def\baselinestretch{1.0}\large\normalsize}{\par}
\pagestyle{empty}
\setlength{\textheight}{9.2in}
\setlength{\textwidth}{6.15in}
\def\basswidth{3.413in}
\def\gutter{0.5in}
\def\noi{\noindent}
\def\eod{\vrule height 6pt width 5pt depth 0pt}
\def\parsec{\par\noindent}
\def\big{\bigskip\parsec}
\def\med{\medskip\parsec}
\def\pr{{\rm Pr}}
\def\l{\ell}
\def\A{{\cal A}}
\def\kk{\kappa}
\def\W{\A}
\def\X{\{X_k\}}
\def\Wk{{\A_{k}}}
\def\eps{\varepsilon}
\def\B{{\cal B}}
\def\TB{\widetilde{\B}}
\def\F{{\cal F}}
\def\S{{\cal S}}
\def\G{{\cal G}}
\def\TG{\widetilde{\G}}
\def\P{{\cal P}}
\def\tx{\hat{x}}
\def\HX{\hat{X}}
\def\tb{\widetilde{B}}
\def\tw{\widetilde{\W}}
\def\TX{\widetilde{X}}
\def\pmin{P_{\min}}
\def\lf{\lfloor}
\def\rf{\rfloor}
\def\lb{L_n^{(b)}}
\def\hb{H_n^{(b)}}
\def\sb{s_n^{(b)}}
\def\O{{\cal O}}
\def\var{{\rm Var~}}
\def\cov{{\rm Cov}}
\def\J{{\cal J}}
\def\I{{\cal I}}
\def\ow{\overline{w}}
\def\ta{{\widetilde{a}}}
\def\real{{I \!\! R}}
\def\h{\hat}
\def\D{{\cal D}}

\vsize=9.05in
\voffset=-0.9in
\hoffset=-0.35in


\begin{document}
\noi
{\bf PATTERN~MATCHING~IMAGE~COMPRESSION~WITH~PREDICATION~LOOP}
\med
%\begin{tabular}{l@{\hspace{2.2in}}l}
\begin{tabular}{ll}
Denis Arnaud&Wojciech Szpankowski\footnotemark\\
%ENST& Department of Computer Science\\
Ecole Nationale Sup\'{e}rieure des T\'{e}l\'{e}communications&Department of Computer Science\\
46 Rue Barrault&Purdue University\\
75013 Paris&W. Lafayette, IN 47907\\
France&U.S.A.\\
%{\tt darnaud@email.enst.fr}& {\tt spa@cs.purdue.edu}
\end{tabular}
\footnotetext{ This work was supported by NSF Grants NCR-9415491
and CCR-9201078.}
\par
\bigskip
Recently, Atallah, Genin and Szpankowski \cite{ags} (cf. also \cite{cs})
introduced a novel image compression technique based on
pattern matching, namely {\it Pattern Matching Image Compression}
(PMIC). Basically, it is a lossy extension of the  well known
Lempel-Ziv scheme in which one searches for the
longest prefix of an uncompressed image that {\it approximately}
occurs in the already processed image.
It was proved
that such an extension leads to a suboptimal compression. 
Success of PMIC crucially depends on several enhancements such as
searching for reverse approximate matching, recognizing substrings in
images that are additively shifted versions of each other, introducing
a variable and adaptive maximum distortion level, and so forth.
Here, we introduce another enhancement, namely, predictive coding.
More importantly, we implement {\it Differential Predictive Code Modulation}
(DPCM) within PMIC. 

More precisely, the PMIC compression algorithm with prediction loop 
should be performed on the quantized differential image on a 
pixel-by-pixel basis. 
Unfortunately, the PMIC algorithm,
proceeds several pixels at a time (i.e., the longest prefix
found in the database). 
But, we can consider a prediction loop algorithm that works
on a cell where the cell in PMIC becomes naturally the longest prefix 
found in the database.

The table below presents a comparison between JPEG, PMIC with constant and
variable distortion level $D$, and PMIC-PL (with prediction loop).
\begin{table}[htbp]
\begin{center}
\begin{tabular}{|l|c|c|c|c|} \hline
``{\bf Lena}''& JPEG & Const D & Var D & PMIC-PL\\ \hline
Compression ratio & 10.03 & 9.99 & 10.01 & 10.04\\ \hline
Bit rate (bits/pixel) & 0.8 & 0.8 & 0.8 & 0.8\\ \hline
RMSE (0-255) & 4.65 & 8.90 & 8.44 & 15.97\\ \hline
PSNR (dB) & 34.78 & 29.14 & 29.61 & 24.06\\ \hline
\end{tabular}
\end{center}
\end{table}
We conclude that prediction gives
some speed up in compression time, better compression ratio  
without significant deterioration of image quality, while a more
sophisticated PMIC-PL has the same time complexity as the original PMIC and
should lead to further improvements in compression. 
However,  the above results
are worse than expected. We believe they originate from
the fact that we use variable-D scheme which seems to be unsuited
for the predictive loop
(the derivative is actually done on the original image, and not
on the differential image). 
For details see \cite{as}.

\footnotesize

\begin{thebibliography}{99}
\bibitem{as}
D. Arnaud and W. Szpankowski,
Pattern Matching Image Compression with Prediction Loop,
%Preliminary Experimental Results,
Purdue University, CSD-TR-96-069, 1996;
also {\tt http://www.cs.purdue.edu/people/spa}.
\bibitem{ags}
M. Atallah, Y. Genin, and W. Szpankowski,
A Pattern Matching Approach to Image Compression,
{\it Proc. International Conference on Image Processing}, vol. II,  349-352,
Lausanne (1996). 
%full version Purdue University, CSD-TR-95-083 (1995) or
%also {\tt http://www.cs.purdue.edu/people/spa}.
\bibitem{cs}
C. Constantinescu and J. Storer,
Improved Techniques for Single-Pass Adaptive Vector
Quantization, {\it Proc. of the IEEE}, 82, 933-939 (1994).
\end{thebibliography}

\end{document}

